\chapter{Introduction}
\label{chap:introduction}

\section{Abstract}
\label{sec:abstract}

\subsection*{Rezumat (Limba Română)}

\nidsname{} este un sistem de detecție a intruziunilor în rețea
(NIDS), construit în jurul a trei subsisteme cooperante: un motor
de reguli scris manual care compară pachetele cu o bază de date
de semnături inițializată din YAML, un clasificator scikit-learn
care scorează aceleași pachete folosind fie Random Forest
(supervizat), fie Isolation Forest (nesupervizat), și un
înregistrator de forensică ce păstrează un mic inel de pachete în
jurul fiecărei alerte. Cele trei subsisteme partajează același
event loop, aceeași reprezentare a pachetelor și același strat
de stocare, astfel încât orice alertă conține contextul complet
care a produs-o.

Codul include un backend FastAPI cu API-uri REST și WebSocket, un
strat asincron SQLAlchemy peste SQLite, un frontend React
(Vite~+~TailwindCSS) și un launcher one-click pentru Windows,
macOS și Linux. O \emph{Consolă de Atac} integrată injectează
șase clase de atac (scan de porturi, brute-force SSH, flood ICMP,
injecție SQL, injecție de comenzi, reverse shell) prin același
dispatcher pe care l-ar alimenta o interfață de rețea reală.
Latența end-to-end, de la sosirea pachetului până la notificarea
în UI, a rămas sub 250\,ms pe un laptop de gamă medie.

\subsection*{Abstract (English)}

\nidsname{} is a Network Intrusion Detection System built around
three cooperating subsystems: a hand-written rule engine that
matches packets against a YAML-seeded signature database, a
scikit-learn classifier that scores those same packets with
either Random Forest (supervised) or Isolation Forest
(unsupervised), and a forensic recorder that persists a small
ring of packets around every alert. The three share one event
loop, one packet representation, and one storage tier, so any
alert carries the full context that produced it.

The codebase covers a FastAPI backend with REST and WebSocket
APIs, an asynchronous SQLAlchemy layer on top of SQLite, a React
frontend (Vite~+~TailwindCSS), and a one-click cross-platform
launcher. An integrated \emph{Attack Console} injects six attack
classes (port scan, SSH brute-force, ICMP flood, SQL injection,
command injection, reverse shell) through the same dispatcher
path a live network interface would feed. End-to-end alert
latency, from packet arrival to UI notification, stayed under
250\,ms on a mid-range laptop.

\subsection*{Cuvinte cheie / Keywords}

\noindent\textbf{Domeniu principal / Main domain:}
Securitatea rețelelor, Sisteme de detecție a intruziunilor (NIDS)
\,/\, Network Security, Network Intrusion Detection.

\noindent\textbf{Domenii adiacente / Adjacent domains:}
Învățare automată, forensică digitală, dezvoltare full-stack web,
inginerie software \,/\, Machine Learning, Digital Forensics,
Full-Stack Web Development, Software Engineering.

\noindent\textbf{Tipul lucrării / Work type:}
Cercetare și sinteză (research and synthesis): the thesis
combines a survey of the NIDS literature with an original
engineering implementation evaluated through a reproducible
synthetic-traffic methodology.

\section{Motivation}
\label{sec:motivation}

There's a gap between what intrusion-detection researchers publish
and what the people guarding actual networks deploy, and it's been
bothering me since I started reading the papers. Pick up the
proceedings of any recent IEEE Security and Privacy conference.
You'll find dozens of papers throwing machine-learning models at
the KDD'99 or CICIDS2017 datasets and reporting F1 scores so high
you'd think the problem was solved~\cite{khraisat2019survey,
buczak2016survey}. Then look at what a real Security Operations
Centre runs in production. Snort signatures. Or Suricata
signatures. Fed into a Splunk dashboard. That's it. The literature
talks about models. The operators want
workflows~\cite{sommer2010outside}.

Why the disconnect? A signature engine emits a clear alert with
a rule name on it. That's what an analyst needs at three in the
morning when a Slack notification wakes them up: a name, a
verdict, something to act on. An ML classifier hands them a
probability score. Probabilities are harder to triage, easier
to tune wrong, and they generate enough false alarms that
operator trust eventually erodes. So when the two paradigms
are framed as alternatives, signatures win on the operational
floor every time --- whether or not they win on paper.

I didn't want to pick a side. The premise of this project is
that signatures and ML don't have to compete. They can
cooperate. Think of it as three categories of alert. The first:
a packet matches a high-confidence rule, end of story. The
second: no rule matched, but the classifier scored the packet
high enough to be suspicious --- the case where a second
opinion is useful. The third: both engines flagged the same
packet. That last one is the strongest signal you'll get. By
labelling every alert with the engine that produced it, the
console tells the analyst not just \emph{what} fired but
\emph{why} they should care.

There was a second reason I started, and it came from looking
at the existing tools. Snort, Suricata and Zeek are good
engines. Really good. But they're all command-line daemons.
There's no analyst console in any of them. The alert page, the
packet pivot, the rule editor --- everything has to be bolted
on by the operator. On the other side of the market sit Splunk
Enterprise Security and IBM QRadar. Beautiful UIs. But the
licences run into five figures a year and the products don't
even do detection themselves --- they aggregate alerts from
external engines. So there's a gap in the middle: a single
self-contained app, self-hostable, that does detection,
alerting and forensics in one binary. \nidsname{} is my attempt
at occupying that gap.

A third reason was more pedagogical. A thesis is also a chance
to show you can ship working software. I wanted a codebase
where somebody could double-click an icon, log in with default
credentials, and see actual alerts in their browser inside
thirty seconds --- no wiki, no package manager, no
infrastructure. That target shaped several decisions. SQLite
instead of Postgres. A virtualenv instead of Docker. A built-in
traffic simulator so the dashboard isn't empty on first boot.
The result is a system that runs end-to-end on a laptop with
nothing else installed. That makes it useful for academic
demonstration, and it works fine as a starting point for
deployment on a small network too.

\section{Problem Statement}
\label{sec:problem-statement}

This thesis tries to answer three concrete engineering questions.

First: \textbf{can a single application combine signature-based and
machine-learning detection in a way that makes the cooperation
visible?} Most existing systems either run one or the other. A few
run both, but in parallel silos with no shared output schema. I
wanted alerts to carry an explicit \emph{detection method} label
--- ``signature'', ``ml'' or ``hybrid'' --- and I wanted the
hybrid label to be earned, not assigned. Both engines have to
flag the same packet for it.

Second: \textbf{can the system keep enough context per alert that
an analyst can investigate without reaching for an external packet
sniffer?} The canonical workflow when an alert fires is ``what
packets led to this?''. In a normal NIDS deployment the operator
opens a separate full-packet-capture tool --- tcpdump, Moloch
(now Arkime), Wireshark --- and correlates by timestamp and
5-tuple. That's two consoles, two indexing layers, and a lot of
chances to get the correlation wrong. I wanted that whole step
removed. Every alert should show up with the last fifty packets
attached, downloadable from the same UI that displayed the alert.

Third: \textbf{can the whole thing be deployed without
infrastructure expertise?} A familiar failure mode of academic
prototypes is that they need a graduate-level setup procedure
before they'll even start. I set myself a hard requirement:
someone who isn't a developer should be able to launch the entire
stack with a single double-click on Windows, macOS or Linux,
without installing dependencies by hand.

\section{Contributions}
\label{sec:contributions}

The contributions of this thesis are:

\begin{enumerate}
  \item A modular detection pipeline (see Chapter~\ref{chap:design},
        Section~\ref{sec:design-detection}) in which signature
        evaluation and ML scoring run on the same packet stream and
        produce three explicit alert categories. The merge logic is
        contained in a single method, \code{Dispatcher.\_evaluate},
        roughly twenty lines long.
  \item A YAML-seeded rule store that is also editable live through
        the UI without restarting the engine. The rule compilation
        step (Section~\ref{sec:design-signature}) builds an in-memory
        list of compiled match predicates which is hot-swapped whenever
        the database changes.
  \item A per-alert JSONL ring-tail capture and a UI page that lets the
        analyst replay the packets and download the file via an
        authenticated endpoint
        (Section~\ref{sec:design-forensics}). The capture is line-
        delimited rather than libpcap-formatted so the analyst can read
        it in any text editor.
  \item A React analyst console
        (Section~\ref{sec:design-frontend}) with eight pages: Login,
        Dashboard, Alerts, Alert Detail, Rules, ML Models, Forensics,
        Attack Console and Settings. Live updates use two WebSocket
        channels (\file{/ws/alerts} and \file{/ws/traffic}) which
        survive page navigation by being scoped to the shared layout
        component.
  \item An integrated \emph{Attack Console}
        (Section~\ref{sec:design-attack-console}) that exposes six
        attack types (Nmap scan, SSH brute-force, ICMP flood, SQL
        injection, command injection, reverse shell) as buttons in the
        UI. Each invocation injects a small burst of synthetic packets
        through the dispatcher exactly as a live capture would, returns
        the number of alerts generated and gives instant audience
        feedback during demonstration.
  \item A one-click cross-platform launcher
        (Section~\ref{sec:design-launcher}) which creates a Python
        virtual environment, installs both the Python and Node
        dependencies, starts FastAPI and Vite as supervised
        subprocesses, waits for both to be reachable, opens the browser
        and handles \verb|Ctrl+C| with an idempotent shutdown routine.
  \item A validation methodology (Chapter~\ref{chap:validation})
        that drives the entire attack catalogue through the
        dispatcher and measures the alerts generated per attack
        kind, the latency from injection to notification, and the
        agreement between the signature and machine-learning
        engines.
\end{enumerate}

\section{Thesis Structure}
\label{sec:thesis-structure}

The rest of the thesis is organised like this.

Chapter~\ref{chap:background} covers the theory I needed: NIDS as
a discipline, signature versus anomaly paradigms, the two ML
algorithms I picked, and the forensic idea that ties everything
together. Chapter~\ref{chap:technologies} walks through the
specific technologies in the stack and why each one ended up
there. Chapter~\ref{chap:related-work} positions the project
against the production-grade tools (Snort, Suricata, Zeek, Wazuh,
Splunk ES) and the published ML-IDS research.
Chapter~\ref{chap:design} is the big technical chapter; it works
through every component of the implementation with file
references and code listings. Chapter~\ref{chap:validation}
describes the validation methodology and reports the latency and
accuracy numbers I collected. Chapter~\ref{chap:conclusions}
wraps up with what got done and what's left for future work.
